\chapter{The question beneath an economy}\label{ch:question}
\question{What are economic systems ultimately meant to make possible?}

The water arrives before the first patient. Someone opened a valve, a motor turned, and a pump lifted water into the clinic's tank. A nurse can wash her hands. Instruments can be cleaned. The building is able to do work that people depend on.

Several records describe the morning. An invoice records electricity purchased. A wage record identifies payment to the person who maintained the pump. A contract assigns responsibility for repairs. A meter records water delivered. An inspection records the condition of a seal. None of these records is foolish or dispensable. Each answers a different question.

Now suppose the electricity invoice is paid on time, but the seal leaks. The payment record may be perfectly correct while the physical service deteriorates. Alternatively, the pump may work perfectly while a patient lacks transport to the clinic. The engineering record can be correct while access fails. Good accounting begins by refusing to turn one kind of success into every kind of success.

\section{Claims, things and services}
An economic claim gives someone a recognised entitlement within an institutional arrangement: to payment, repayment, delivery or use. A physical resource is something represented in the material world, such as water in a tank. A service is what an arrangement makes possible, such as safe washing at a particular time. A human need explains why that service matters.

These descriptions connect, but do not replace one another. A contract for water can help organise a reliable supply. It cannot itself hydrate a patient. A tank of water may help provide care, but water of the wrong quality, in the wrong place, or available after the emergency is not the same service. A litre is not a promise, and a promise is not a litre.

The distinction is especially important for maintenance. A replacement bearing can be physically small and functionally decisive. The account must say what it represents before assigning a number. Does it record the bearing's mass, the motor's operating condition, the electricity consumed during repair, or the increased ability to deliver water later? Those are different objects, with different evidence requirements.

\illustration{two_records}{Schematic: one pumping event supports several records. The payment and physical records can be linked by event identity without being interchangeable. No EBU valuation is assigned by this diagram.}{fig:records}

\section{What EBU proposes to ask}
The Energy Balance Project investigates an explicit account of physical action. In this book, an Energy Balance Unit, or EBU, expresses a signed change in a declared measure of represented burden, after any separately justified process burden. The declaration matters: an account has to identify what it measures and what it leaves outside.

Before we can calculate, we must answer ordinary questions. What was present before the action? What did the action change? Which physical restrictions applied? Relative to what reference is the change evaluated? Has an important effect already been counted elsewhere? Who chose the action? Who verified that it occurred?

The resulting number cannot contain all possible answers about the clinic. It need not be a price. It does not state the moral worth of the nurse, mechanic or patient. It cannot certify an effect that the chosen state description does not represent. A local equation is useful only if its limited meaning stays visible.

\section{The ethical starting point}
The motivation of this book includes an ethical commitment: people should have reliable access to essential support, and avoidable damage to the conditions of life deserves attention. That is a declared value commitment. It is not a fact proved by a mathematical curve.

Different communities may disagree about priorities, responsibilities and acceptable trade-offs. A transparent physical account could inform those disagreements. It cannot eliminate them by renaming a disputed choice a measurement. The right to receive care must not be silently made conditional on a person's ability to generate a positive score in a model.

The central research argument is therefore precise. Monetary transactions and existing institutions do not automatically guarantee preservation of all specified life-supporting functions. EBU investigates an explicit physical-action account; its usefulness remains a research question. Evidence that a problem exists is a reason to investigate. It is not evidence that this proposal solves it.

\practice{The clinic paid for ten units of water, but its meter recorded eight. Which record settles the physical delivery question?}{The meter is relevant to the delivered quantity, subject to its accuracy and placement. The invoice establishes what was charged or contracted. Comparing the records may reveal a discrepancy, but neither alone tells us whether two units leaked, were stored elsewhere, or were mismeasured. A missing explanation should remain visible.}

\carry{We need a physical account because physical functioning is one of the things we care about. To understand its proposed role, we should first understand the considerable work monetary signals already do.}
